% NOTE: Filename is fixed by the book outline; content teaches Week 11
% (Change Data Capture with AWS DMS & EventBridge, and Data Contracts).
\chapter{Change Data Capture, Event-Driven Pipelines, and Data Contracts}
\index{AWS DMS}\index{change data capture}\index{Amazon EventBridge}\index{AWS Lambda}\index{Glue Schema Registry}\index{data contracts}\index{schema evolution}%
\begin{objectives}
\begin{itemize}
    \item Configure AWS DMS change data capture (CDC) replication from a relational source to an S3 data lake.
    \item Monitor CDC health with the right CloudWatch metrics and detect silent lag.
    \item Build event-driven pipelines with Amazon EventBridge and lightweight AWS Lambda transformations.
    \item Define data contracts with the AWS Glue Schema Registry: Avro schemas, compatibility modes, and versioning.
    \item Distinguish breaking from non-breaking schema changes and assign contract ownership.
    \item Architect near-real-time synchronization from an on-premises Oracle database to Redshift for live BI.
\end{itemize}
\end{objectives}

\begin{crosscloud}
CDC, event routing, and schema contracts are cloud-agnostic disciplines; each platform offers a managed shape of the same ideas.

\vspace{2mm}
{\footnotesize
\begin{tabular}{@{}p{2.8cm}p{3.0cm}p{3.0cm}p{3.0cm}@{}}
\toprule
\textbf{Capability} & \textbf{AWS} & \textbf{Azure} & \textbf{GCP} \\
\midrule
Managed CDC replication & AWS DMS & Azure DMS / Data Factory CDC & Datastream \\
Event router & EventBridge & Event Grid & Eventarc \\
Serverless functions & Lambda & Azure Functions & Cloud Run functions \\
Schema registry & Glue Schema Registry & Azure Schema Registry & Managed Service for Schema Registry \\
Contract enforcement & In Glue jobs / registry checks & In pipelines & In Dataflow / registry checks \\
\addlinespace
CDC lag metric & CDCLatencySource/Target & Replication latency metrics & Datastream freshness \\
Breaking-change policy & BACKWARD/FORWARD/FULL modes & Compatibility levels & Compatibility levels \\
\bottomrule
\end{tabular}}

\vspace{1mm}
\textbf{Snowflake} ingests CDC through Snowpipe and streams, applying the same contract ideas at the table boundary. \textbf{MongoDB} contributes change streams---the document-native CDC primitive that pairs naturally with everything in this chapter. \textbf{Fabric} routes database mirroring and event streams through OneLake, its version of the CDC-to-lake path.
\end{crosscloud}

\section{Week 11 Roadmap and Mental Model}
Batch windows are shrinking. The business question behind this week is: how fresh can analytical data be, and what breaks when it is continuous? The answer has two halves. The \emph{plumbing} half is change data capture---reading the source database's transaction log and replaying inserts, updates, and deletes into the lake continuously. The \emph{contract} half is making sure the continuous stream cannot silently change shape: schemas are registered, compatibility is enforced, and producers cannot break consumers without an explicit versioned decision.

Weeks 1--10 mostly moved data on a schedule; this week moves it on \emph{change}. The cadence changes the failure modes: instead of a failed nightly job, you get a silently lagging stream and a subtly drifting schema. The week's tooling---DMS, EventBridge, Lambda, Schema Registry---maps onto those two failure modes directly.
\index{CDC}\index{data contract}

\begin{definitionbox}
\textbf{Change data capture (CDC)} is the technique of reading a source database's transaction log to extract insert, update, and delete events as they happen, and applying them to a target---keeping the target continuously synchronized without re-scanning the source tables.
\end{definitionbox}

\begin{keyterms}
\textbf{Full load + CDC}: DMS task pattern that bulk-loads existing rows, then applies ongoing changes. \textbf{Replication instance}: the managed DMS compute that runs tasks. \textbf{CDCLatencySource / CDCLatencyTarget}: CloudWatch metrics measuring lag at the read and apply stages. \textbf{Event bus}: EventBridge's routing fabric for events with pattern-matched rules. \textbf{Schema registry}: versioned store of schemas with enforced compatibility. \textbf{Compatibility mode}: BACKWARD, FORWARD, or FULL---which evolutions are permitted.
\end{keyterms}

\section{Monday: AWS DMS Architecture for Continuous Replication}
\subsection{Full Load Plus CDC}
A DMS task has two phases. \emph{Full load} copies existing table contents; \emph{CDC} then tails the transaction log and applies changes as they arrive \cite{awsDMSUserGuide}. The pair matters separately: full load is a one-time bulk problem (parallelism, LOB handling, validation), while CDC is a forever problem (lag, log retention, source load). Tasks run on a replication instance you size like any compute: memory for cached changes during bursts, storage for spilled transactions.
\index{full load}\index{replication instance}

\subsection{CDC to the Lake}
Replicating to S3, DMS writes full-load files and ongoing change files (with operation markers) under a partitioned prefix---commonly \texttt{table=/load-date=/}. Downstream, a Glue job (Chapter~10) or a MERGE process applies the changes to a current-state table. This two-layer design---immutable change log in the lake, current-state table materialized downstream---preserves history, makes replays possible, and isolates consumers from replication mechanics.
\index{DMS to S3}

\begin{pitfall}
\textbf{Monitoring task status instead of lag.} A DMS task shows \texttt{running} while minutes---or hours---behind the source. Alarm on \texttt{CDCLatencySource} and \texttt{CDCLatencyTarget}, and validate that the source database retains enough transaction log history to survive a replication instance outage. Stale CDC with a ``running'' task is the classic silent failure.
\end{pitfall}

\section{Tuesday: Event-Driven Pipelines with EventBridge}
\subsection{Routing Events Instead of Polling}
Amazon EventBridge is a serverless event bus: producers publish events, rules match patterns (source, detail-type, fields in the payload), and matched events fan out to targets---Lambda, Step Functions, SQS, Kinesis, and more \cite{awsEventBridgeDocs}. For data engineering, its signature use is reacting to \emph{arrival}: an S3 \texttt{ObjectCreated} event matching \texttt{raw/invoices/*.csv} starts the pipeline that validates and loads that file, replacing polling loops and fixed schedules with demand-driven execution.
\index{EventBridge}\index{event bus}

\subsection{Events as Pipeline Contracts}
An event has a schema too. EventBridge schema discovery can infer and register event shapes, and mature teams treat event payloads like any other contract: versioned, documented, and owned by the producer. A rule pattern is a subscription to a contract; breaking the payload breaks every subscriber's rule silently. That observation leads directly to Wednesday's topic.

\section{Wednesday: Lambda Transformations and Data Contracts}
\subsection{Lambda in the Data Path}
AWS Lambda handles the transforms that are too small for a Spark cluster: validating a header row, enriching an event, splitting a batch file, calling an API per record \cite{awsLambdaDocs}. The discipline is knowing the boundary: Lambda bills per millisecond with a 15-minute ceiling---perfect for per-file or per-event work, wrong for the 2~TB join. In this week's architecture Lambda is the connective tissue: triggered by EventBridge, normalizing DMS output, enforcing schemas at ingestion.
\index{Lambda}

\subsection{Data Contracts with the Glue Schema Registry}
A \emph{data contract} is an enforceable agreement about a stream's or table's shape: field names, types, optionality, and evolution rules. The AWS Glue Schema Registry implements the machinery: Avro (or JSON Schema) definitions, versioned, with a \emph{compatibility mode} that decides which evolutions are legal \cite{awsGlueSchemaRegistry}:
\index{Glue Schema Registry}\index{compatibility mode}

\begin{itemize}
    \item \textbf{BACKWARD}: new schema can read old data (consumers upgrade first). Deleting a field is allowed; adding a field without a default is not.
    \item \textbf{FORWARD}: old schema can read new data (producers upgrade first). Adding an optional field is allowed; deleting is not.
    \item \textbf{FULL}: both directions---the conservative production default. Only changes safe both ways (e.g., adding an optional field with a default) pass.
\end{itemize}

The contract taxonomy is worth memorizing: \emph{adding an optional field} is non-breaking; \emph{deleting a field, renaming a field, or changing a field's type} breaks consumers. Ownership is the social half: the producer owns the contract, but consumers sign off on changes---the registry enforces the mechanics, the sign-off process enforces the intent.

\begin{lstlisting}[language=bash,caption={Registering an Avro schema with FULL compatibility.},label={lst:ch11-registry}]
aws glue create-registry --registry-name orders-registry

aws glue create-schema `
  --registry-id RegistryName=orders-registry `
  --schema-name order-events `
  --data-format AVRO `
  --compatibility FULL `
  --schema-definition "{`"type`":`"record`",`"name`":`"Order`",`"fields`":[`
      {`"name`":`"order_id`",`"type`":`"string`"},`
      {`"name`":`"amount`",`"type`":`"double`"},`
      {`"name`":`"coupon`",`"type`":[`"null`",`"string`"],`"default`":null}`
    ]}"
\end{lstlisting}

\begin{tipbox}
Enforce contracts at \emph{ingestion}, not at query time. A Glue job or Lambda that validates each batch against the registry rejects a breaking change before it poisons the lake; discovering the break at query time means the bad data is already persisted, versioned, and possibly replicated.
\end{tipbox}

\section{Thursday Hands-on Project: CDC from RDS PostgreSQL to the Lake}
\index{hands-on project!DMS CDC to S3}
Configure a DMS task with CDC enabled to stream ongoing changes from an RDS PostgreSQL instance into an S3 data-lake bucket, then prove the stream is live and lag is observable.

\subsection*{Lab Objectives}
\begin{itemize}
    \item Prepare the PostgreSQL source for logical replication.
    \item Create a replication instance, source and target endpoints, and a full-load-plus-CDC task.
    \item Generate changes on the source and observe them arriving as change files in S3.
    \item Alarm on CDC lag metrics.
\end{itemize}

\subsection*{Procedure}
\begin{enumerate}
    \item \textbf{Enable logical replication} on the RDS instance: parameter group with \texttt{rds.logical\_replication = 1}, reboot, and confirm \texttt{wal\_level = logical}.
    \item \textbf{Create the replication instance} (a \texttt{dms.t3.medium} suffices for the lab) in the same VPC as the database.
    \item \textbf{Create endpoints}: a PostgreSQL source endpoint (with the master or a replication user) and an S3 target endpoint pointing at \texttt{s3://<bucket>/cdc/orders/}.
    \item \textbf{Create the task} with migration type \texttt{full-load-and-cdc} and a table mapping over the \texttt{public.orders} table.
\end{enumerate}

\begin{lstlisting}[language=bash,caption={Creating the DMS task with full load plus CDC.},label={lst:ch11-dms-task}]
$Region = "us-east-1"

aws dms create-replication-instance `
  --replication-instance-identifier week11-ri `
  --replication-instance-class dms.t3.medium `
  --allocated-storage 50

# (Endpoints omitted for brevity: source = RDS PostgreSQL, target = S3 prefix.)

$Mapping = @"
{
  "rules": [{
    "rule-type": "selection",
    "rule-id": "1",
    "rule-name": "orders",
    "object-locator": { "schema-name": "public", "table-name": "orders" },
    "rule-action": "include"
  }]
}
"@
$Mapping | Set-Content -Path table-mappings.json -Encoding ascii

aws dms create-replication-task `
  --replication-task-identifier week11-orders-cdc `
  --source-endpoint-arn "arn:aws:dms:${Region}:123456789012:endpoint:SRC" `
  --target-endpoint-arn "arn:aws:dms:${Region}:123456789012:endpoint:TGT" `
  --replication-instance-arn "arn:aws:dms:${Region}:123456789012:rep:RI" `
  --migration-type full-load-and-cdc `
  --table-mappings file://table-mappings.json

aws dms start-replication-task `
  --replication-task-arn "arn:aws:dms:${Region}:123456789012:task:TASK" `
  --start-replication-task-type start-replication
\end{lstlisting}

Generate change traffic and verify:

\begin{lstlisting}[language=SQL,caption={Source-side changes that should appear as CDC files.},label={lst:ch11-changes}]
INSERT INTO orders (order_id, amount, status) VALUES (90001, 129.99, 'NEW');
UPDATE orders SET status = 'PAID' WHERE order_id = 90001;
DELETE FROM orders WHERE order_id = 90001;
\end{lstlisting}

\subsection*{Verification}
\begin{itemize}
    \item Full-load files appear under the S3 prefix before CDC files start.
    \item Each source change produces a change record with its operation marker (I/U/D) in the target prefix within the observed latency.
    \item CloudWatch shows \texttt{CDCLatencySource} and \texttt{CDCLatencyTarget} near zero under light load; an alarm exists on both.
    \item Stopping writes and querying the latest change files shows the delete as the final state of \texttt{order\_id 90001}.
\end{itemize}

\section{Friday Use Case: Near-Real-Time Oracle-to-Redshift Synchronization}
\index{use case!Oracle to Redshift CDC}
A bank runs its core ledger on on-premises Oracle. The BI team needs Redshift dashboards reflecting ledger changes within five minutes; the nightly batch is no longer acceptable, and the Oracle estate cannot be modified beyond enabling supplemental logging.

\subsection{Architecture}
\begin{figure}[H]
    \centering
    \resizebox{\textwidth}{!}{%
    \begin{tikzpicture}[
        box/.style={rectangle, rounded corners, draw=primary, thick, fill=primary!6, align=center, minimum width=2.9cm, minimum height=1.0cm, font=\small\bfseries},
        storebox/.style={cylinder, shape aspect=0.25, draw=primary, thick, fill=blue!6, shape border rotate=90, align=center, text width=2.6cm, minimum height=1.3cm, font=\small\bfseries},
        guard/.style={rectangle, rounded corners, draw=red!60!black, thick, fill=red!5, align=center, minimum width=2.6cm, minimum height=0.9cm, font=\footnotesize\bfseries},
        arr/.style={-{Stealth[length=2.5mm]}, draw=primary, thick}
    ]
        \node[box] (oracle) at (0,0) {On-prem Oracle\\(redo logs)};
        \node[box] (dms) at (3.3,0) {DMS CDC\\replication instance};
        \node[storebox] (s3) at (6.6,0) {S3 change log\\(raw zone)};
        \node[box] (glue) at (9.9,0) {Glue apply\\(MERGE)};
        \node[box] (rs) at (13.0,0) {Redshift\\current state};
        \node[guard] (lag) at (6.6,2.0) {Lag alarms\\CDCLatency*};
        \node[guard] (reg) at (9.9,2.0) {Schema Registry\\contract checks};
        \draw[arr] (oracle) -- (dms);
        \draw[arr] (dms) -- (s3);
        \draw[arr] (s3) -- (glue);
        \draw[arr] (glue) -- (rs);
        \draw[arr, dashed] (lag) -- (dms);
        \draw[arr, dashed] (reg) -- (glue);
    \end{tikzpicture}%
    }
    \caption{Near-real-time ledger sync: DMS tails Oracle redo into an S3 change log; a Glue MERGE maintains current state in Redshift; lag and contracts are monitored.}
    \label{fig:ch11-oracle-cdc}
\end{figure}

\subsection{Design Decisions}
\textbf{Two targets, not one.} DMS writes the raw change log to S3 (durable, replayable, auditable) and a second task---or a downstream Glue job reading the change log---applies MERGEs into Redshift. Writing straight to Redshift couples the replication path to warehouse availability; the S3 stage decouples them and preserves the full history.

\textbf{Freshness budget.} The five-minute SLA is spent deliberately: DMS apply latency (sub-minute under normal load), S3 landing, a scheduled or event-driven Glue MERGE every two minutes, and Redshift commit time. Each stage gets a metric; the sum gets the alarm.

\textbf{Log retention as an availability dependency.} Oracle must retain enough redo/archive logs to cover the longest plausible replication outage. The runbook names the retention, the metric that watches it, and the re-snapshot procedure if it is ever exceeded.

\textbf{Schema drift policy.} Adding a nullable column flows: FULL compatibility permits it, the Glue apply adopts it, Redshift gains it via an expand-first migration (Chapter~24 formalizes this pattern). A type change or rename halts the pipeline and pages the owner---it is a contract violation, not a runtime event.

\begin{pitfall}
\textbf{Treating the change log as disposable.} The S3 change log is your replay capability and your audit trail. Lifecycle-expire it early and a Redshift-side corruption becomes unrecoverable without a new full load---which on a core ledger means an outage negotiation, not a retry.
\end{pitfall}

\section{Interview Q\&A}
\subsection{Concepts and Trade-offs}
\begin{enumerate}
    \item \textbf{Q: Which CloudWatch metrics reveal CDC lag while a task shows ``running''?}\\
    A: \texttt{CDCLatencySource} (gap between the source commit and DMS reading it) and \texttt{CDCLatencyTarget} (gap between reading and applying at the target). Task status alone says nothing about freshness.
    \item \textbf{Q: BACKWARD versus FULL compatibility: which permits deleting a field safely?}\\
    A: BACKWARD (new schema reads old data, so consumers already tolerate the absence). FULL forbids deletions because old consumers could not read new data missing the field.
    \item \textbf{Q: Who owns a data contract: producer, consumer, or both?}\\
    A: The producer owns and versions the schema; consumers sign off on changes. The registry enforces mechanics; the sign-off workflow enforces intent.
    \item \textbf{Q: Why enforce the schema at ingestion instead of at query time?}\\
    A: By query time the bad data is already persisted---and possibly versioned, replicated, and transformed downstream. Ingestion enforcement rejects breaking changes before they become durable.
    \item \textbf{Q: Name one breaking and one non-breaking schema change.}\\
    A: Breaking: deleting a field, renaming it, or changing its type. Non-breaking: adding an optional field with a default.
    \item \textbf{Q: Why stage DMS output in S3 before Redshift?}\\
    A: The change log is replayable history and decouples replication from warehouse availability; MERGEs can be re-run, audited, and repointed without touching the source again.
\end{enumerate}

\section{Further Reading}
The DMS User Guide and Best Practices cover replication instances, full-load-plus-CDC tasks, and source prerequisites \cite{awsDMSUserGuide,awsDMSBestPractices}. EventBridge and Lambda documentation cover event routing patterns and serverless transform limits \cite{awsEventBridgeDocs,awsLambdaDocs}. The Glue Schema Registry documentation details compatibility modes and versioning \cite{awsGlueSchemaRegistry}. OpenLineage provides the emerging standard for tracing where contract-violating data flowed \cite{openLineage}.

\section*{Key Takeaways}
\begin{itemize}
    \item CDC replays the transaction log: continuous freshness without re-scanning sources---and a new silent-failure mode called lag.
    \item Monitor \texttt{CDCLatencySource}/\texttt{CDCLatencyTarget}, and treat source log retention as an availability dependency.
    \item EventBridge replaces polling with routed arrival events; Lambda handles the per-event transforms too small for clusters.
    \item Data contracts make streams safe to evolve: registry, versions, and compatibility modes, owned by producers with consumer sign-off.
    \item Enforce contracts at ingestion; query-time discovery means the poison is already durable.
    \item Stage CDC through an S3 change log; current state is a rebuildable materialization, history is the asset.
\end{itemize}

\section{Exercises}
\begin{enumerate}
    \item \textbf{(Conceptual)} Explain the difference between full load and CDC phases, and why each has distinct failure modes.
    \item \textbf{(Conceptual)} Why does FULL compatibility permit adding an optional field but forbid deleting one?
    \item \textbf{(Hands-on)} Complete the Thursday lab, then write a Glue job that MERGEs the change files into a current-state table; verify an update and a delete propagate.
    \item \textbf{(Hands-on)} Register the \texttt{order-events} schema with FULL compatibility; attempt to register a version that deletes \texttt{amount} and capture the rejection.
    \item \textbf{(Hands-on)} Build the EventBridge rule from arrival to action: S3 \texttt{ObjectCreated} on a prefix triggers a validating Lambda, with failures to an SQS DLQ.
    \item \textbf{(Design)} Design the freshness monitoring for the Friday architecture: metrics per stage, alarm thresholds, and the five-minute budget arithmetic.
    \item \textbf{(Design)} A producer must rename \texttt{amount} to \texttt{total\_amount}. Write the multi-release contract-evolution plan that never breaks consumers.
    \item \textbf{(Architecture)} Add OpenLineage emission to the Glue apply job and describe how it shortens the investigation of a contract violation.
\end{enumerate}

\section{Capstone Project: A Contract-Protected CDC Feed}\label{sec:ch11-capstone}
\index{capstone project}\index{CDC!capstone}\index{data contracts!capstone}

\begin{capstonebox}
\textbf{Mission.} Stand up a production-shaped CDC feed: DMS full load plus CDC from PostgreSQL into an S3 change log, a contract-validating Glue apply into Redshift, EventBridge-driven arrival processing, and lag monitoring with alarms.

\textbf{Estimated time.} 5--7 hours in a sandbox AWS account.

\textbf{What you\textquotesingle{}ll practice.}
\begin{itemize}
    \item DMS replication design: instance sizing, endpoints, full-load-plus-CDC tasks.
    \item Contract enforcement at ingestion with the Glue Schema Registry.
    \item Event-driven apply pipelines with EventBridge and Lambda.
    \item Freshness observability: lag metrics, alarms, and re-snapshot runbooks.
\end{itemize}
\end{capstonebox}

\subsection*{Scenario}
An order-management PostgreSQL database must feed Redshift analytics within ten minutes, continuously, with a guarantee that a schema change upstream can never silently corrupt the warehouse. The feed must survive a one-hour replication outage without data loss.

\subsection*{Requirements}
\textbf{Functional requirements.}
\begin{itemize}
    \item DMS full-load-plus-CDC task from RDS PostgreSQL to an S3 change log, partitioned by table and date.
    \item A Glue apply job that validates each batch against a FULL-compatibility registry schema before MERGEing into Redshift.
    \item EventBridge rules triggering validation on change-file arrival; failures route to a DLQ with an SNS alert.
    \item CloudWatch alarms on both CDC lag metrics and on apply-job failure.
\end{itemize}

\textbf{Non-functional requirements.}
\begin{itemize}
    \item The feed tolerates a one-hour DMS outage with zero data loss (documented log-retention evidence).
    \item A deliberately breaking schema change is rejected at ingestion with the rejection logged.
    \item Rerunning the apply for any date is idempotent.
\end{itemize}

\subsection*{Implementation Steps}
\begin{enumerate}
    \item Prepare the source (logical replication) and create the DMS task per \cref{lst:ch11-dms-task}.
    \item Register the contract per \cref{lst:ch11-registry} and encode validation in the apply job.
    \item Build the MERGE apply and prove idempotency by double application.
    \item Wire EventBridge arrival events, the DLQ, and the SNS alerts.
    \item Configure lag alarms; simulate an outage by stopping the task and document recovery.
    \item Write the runbook: lag response, contract-violation response, re-snapshot procedure.
\end{enumerate}

\subsection*{Deliverables}
\begin{itemize}
    \item Task configuration, table mappings, and endpoint definitions as code or exported JSON.
    \item The registry schema, its compatibility policy, and the validation code path.
    \item Evidence: change files, MERGE results, an injected breaking change rejected, alarm history from the outage drill.
    \item The freshness budget arithmetic and the runbook.
\end{itemize}

\subsection*{Acceptance Criteria}
\begin{itemize}
    \item End-to-end freshness meets the ten-minute target under normal load, measured.
    \item The one-hour outage drill recovers with no gap and no duplicate application.
    \item A breaking schema change cannot reach Redshift; the rejection is logged and alerted.
    \item Every stage of the freshness budget has a metric and an alarm.
\end{itemize}

\subsection*{Stretch Goals}
\begin{itemize}
    \item Add a second target task writing Parquet current-state snapshots for Athena access.
    \item Implement consumer sign-off as a pull-request workflow on schema versions.
    \item Replace scheduled apply with Kinesis Data Streams delivery (Chapter~13) and compare latency.
    \item Emit OpenLineage events from the apply job and trace one record end to end.
\end{itemize}
